% ════════════════════════════════════════════════════════════════════
%  front_pages/abstract.tex
%  Abstract Page
% ════════════════════════════════════════════════════════════════════

\chapter*{Abstract}
\addcontentsline{toc}{chapter}{Abstract}
\thispagestyle{plain}

{\singlespacing
\noindent
Non-small cell lung cancer (NSCLC) accounts for about 85\% of lung cancer cases, and
EGFR mutations are among its most significant oncogenic drivers. EGFR-targeted tyrosine
kinase inhibitors have transformed treatment, yet acquired resistance and high therapy
cost remain major challenges, especially in resource-limited South Asian settings. This
study presents a graph neural network (GNN) framework for EGFR-inhibitor classification
and the prospective virtual screening of compounds from traditional Nepali medicinal
plants.

\vspace{0.18cm}
\noindent
A curated dataset of 10,523 compounds with validated EGFR bioactivity (74\% active) was
assembled from ChEMBL, standardised with RDKit, and represented as molecular graphs with
29-dimensional atom and 12-dimensional bond features. Seven models were benchmarked under
a Bemis-Murcko scaffold split: five GNNs (GCN, GAT, GIN, AttentiveFP, GraphSAGE) and two
ECFP4 baselines (Random Forest, MLP). The baselines led (Random Forest AUROC 0.8827, MLP
0.8547) ahead of the GNNs (0.8087--0.8454, GCN strongest); Optuna tuning raised
AttentiveFP's validation AUROC to 0.8632 but did not surpass them---well-tuned fingerprint
models remain highly competitive at this dataset size.

\vspace{0.18cm}
\noindent
The two best models (GCN and Random Forest) were applied to a natural-product library of
seven Nepali plant species from COCONUT~2.0. Of 535 retrieved compounds, eight already in
the ChEMBL training data---including the apparent top hits apigenin and kaempferol---were
excluded as leakage, leaving 527 genuinely unseen compounds. The hit list was strongly
model-dependent (Random Forest 179 predicted actives versus GCN 24), and neither best
model gave a Lipinski-compliant lead at probability $\geq 0.80$. The strongest candidates
are the xanthone \textbf{isogentisin} (Random Forest 0.728) and the flavone
\textbf{luteolin}, the top cross-model consensus lead---both at moderate confidence.

\vspace{0.18cm}
\noindent
Model-consensus screening can thus surface drug-like natural-product candidates from
Nepali ethnomedicinal libraries using an open-source, CPU-only pipeline, while showing
that such screens are model-dependent, must exclude training-set compounds to stay
prospective, and yield only moderate-confidence leads requiring experimental validation.

\vspace{0.3cm}
\noindent{\fontsize{10}{12}\selectfont\textit{\textbf{Keywords:}
EGFR inhibitors, graph neural networks, virtual screening,
Nepali medicinal plants, drug discovery.}}
\par}